\chapter{Equivalent formulations of quantum theory}

We will now show equivalent formulations of quantum theory.

\section{Dirac's vector notation}

Comparing Schrödinger's wave mechanics with Dirac's abstract vector formalism:

\begin{center}
\begin{tabular}{|l|l|l|}
\hline
Schrödinger &  & Dirac \\
\hline
Wavefunctions $\Psi_{n}$ & $\Longleftrightarrow$ & Vectors $|n\rangle$ \\
\hline
Fourier series decomposition: & $\Longleftrightarrow$ & Vector components decomposition: \\
\hline
$\Psi_{v}=\sum_{i} v_{i} \phi_{i}$ &  & $|v\rangle=\sum_{i} v_{i} \vec{i}$ \\
\hline
Operators $\hat{A}$ & $\Longleftrightarrow$ & Matrices $\tilde{A}$ \\
\hline
Probability density $\rho=|\Psi|^{2}$ &  & ? \\
\hline
\end{tabular}
\end{center}

The probabilistic interpretation remains intact in Dirac's formalism through the inner product:

\begin{equation}
\left(\Psi_{u}, \Psi_{v}\right)=\langle u \mid v\rangle=\sum_{i} u_{i}^{*} v_{i}
\end{equation}

For the special case where $u=v$:

\begin{equation}
\left(\Psi_{v}, \Psi_{v}\right)=\langle v \mid v\rangle=\sum_{i}\left|v_{i}\right|^{2}
\end{equation}

Here $\left|v_{i}\right|^{2}$ represents the probability of finding the system in state $\phi_{i}$ or $|i\rangle$.

Operators in Schrödinger's formulation transform into matrices in Dirac's approach. This transformation occurs through the basis representation:

\begin{equation}
\mathcal{B}=\left\{\phi_{i}: i \in[0,+\infty]\right\}
\end{equation}

The matrix elements are defined as:

\begin{equation}
A_{n m}=\left(\phi_{n}, A \phi_{m}\right)
\end{equation}

Leading to the complete operator representation:

\begin{equation}
\tilde{A}=\sum_{i} \sum_{j}|i\rangle A_{i j}\langle j|
\end{equation}


Recall that $|i\rangle\langle j|$ generates a $n \times n$ matrix with all zeros and 1 in position $i j$. For example let $i=1, j=2$ with $n=3$ :

\[
|1\rangle\langle 2|=\left[\begin{array}{l}
1  \\
0 \\
0
\end{array}\right]\left[\begin{array}{lll}
0 & 1 & 0
\end{array}\right]=\left[\begin{array}{lll}
0 & 1 & 0 \\
0 & 0 & 0 \\
0 & 0 & 0
\end{array}\right]
\]

Using bra-ket notation, we can derive the matrix elements:

\begin{equation}
\langle n| \tilde{A}|m\rangle=\langle n|\left(\sum_{i} \sum_{j}|i\rangle A_{i j}\langle j|\right)|m\rangle=\sum_{i} \sum_{j} \underbrace{\langle n \mid i\rangle}_{\delta_{n i}} A_{i j} \underbrace{\langle j \mid m\rangle}_{\delta_{j m}}=\sum_{i} \sum_{j} \delta_{n i} A_{i j} \delta_{j m}=A_{n m}
\end{equation}

The adjoint matrix follows naturally from the Schrödinger picture definition:

\begin{equation}
A_{m n}^{\dagger}=\left(\phi_{m}, A^{\dagger} \phi_{n}\right)=\left(A \phi_{m}, \phi_{n}\right)=\left(\phi_{n}, A \phi_{m}\right)^{*}=A_{n m}^{*}
\end{equation}

Thus the adjoint matrix equals the transpose complex conjugate. For Hermitian operators where $A^{\dagger}=A$, we have $A_{n m}^{\dagger}=A_{n m}$.

When the basis $\mathcal{B}$ diagonalizes operator $A$, the matrix representation simplifies to:

\begin{equation}
\tilde{A}=\sum_{n} \sum_{m}|n\rangle a_{n} \delta_{n m}\langle m|=\sum_{n} a_{n}|n\rangle\langle n|
\end{equation}

Where $a_n$ are the eigenvalues satisfying:

\begin{equation}
A \phi_{n}=a_{n} \phi_{n}
\end{equation}

\section{Basis transformations}
When working with a different basis $\mathcal{B}^{\prime}=\left\{\phi_{k}^{\prime}: k \in[0,+\infty]\right\}$, we need transformation rules between bases.

Lemma 10.1. Two complete orthonormal bases $\mathcal{B}$ and $\mathcal{B}^{\prime}$ are connected by a unitary transformation.

Proof. Any element of $\mathcal{B}^{\prime}$ can be expanded in terms of $\mathcal{B}$:

\begin{equation}
\phi_{n}^{\prime}=\sum_{m} S_{m n} \phi_{m}
\end{equation}

Where the transformation coefficients are:

\begin{equation}
S_{m n}=\left(\phi_{m}, \phi_{n}^{\prime}\right)
\end{equation}

For $S$ to be unitary, it must satisfy:

\begin{equation}
S S^{\dagger}=S^{\dagger} S=\mathbb{I} \Longleftrightarrow \sum_{i} S_{m i} S_{i n}^{\dagger}=\sum_{i} S_{m i}^{\dagger} S_{i n}=\delta_{n m}
\end{equation}

Using $S_{i n}^{\dagger}=S_{n i}^{*}$ and the completeness of basis $\mathcal{B}^{\prime}$:

\begin{align}
\sum_{i} S_{m i} S_{i n}^{\dagger} & =\sum_{i} S_{m i} S_{n i}^{*}=\sum_{i}\left(\phi_{m}, \phi_{i}^{\prime}\right)\left(\phi_{n}, \phi_{i}^{\prime}\right)^{*}=\sum_{i}\left(\phi_{m}, \phi_{i}^{\prime}\right)\left(\phi_{i}^{\prime}, \phi_{n}\right)= \\
& =\sum_{i}\left(\int \mathrm{d} x \phi_{m}^{*}(x) \phi_{i}^{\prime}(x)\right)\left(\int \mathrm{d} y \phi_{i}^{*}(y) \phi_{n}(y)\right)= \\
& =\int \mathrm{d} x \int \mathrm{d} y \phi_{m}^{*}(x) \phi_{n}(y) \underbrace{\sum_{i}\left[\phi_{i}^{\prime}(x) \phi_{i}^{\prime}(y)\right]}_{\delta(x-y)}= \\
& =\int \mathrm{d} x \int \mathrm{d} y \phi_{m}^{*}(x) \phi_{n}(y) \delta(x-y)= \\
& =\int \mathrm{d} x \phi_{m}^{*}(x) \phi_{n}(x)=\left(\phi_{m}, \phi_{n}\right)=\delta_{m n}
\end{align}

This completes the proof.

Corollary 10.2. Given a unitary transformation $S$ from basis $\mathcal{B}^{\prime}$ to $\mathcal{B}$, the inverse transformation is $S^{\dagger}$.

Proof.

\begin{equation}
\sum_{n} S_{n i}^{\dagger} \phi_{n}^{\prime}=\sum_{n} S_{n i}^{\dagger} \sum_{m} S_{m n} \phi_{m}=\sum_{m} \underbrace{\left(\sum_{n} S_{m n} S_{n i}^{\dagger}\right)}_{\delta_{m i}} \phi_{m}=\sum_{m} \delta_{m i} \phi_{m}=\phi_{i}
\end{equation}


So:

\begin{equation}
\phi_{i}=\sum_{n} S_{n i}^{\dagger} \phi_{n}^{\prime}
\end{equation}

The transformation also relates vector components between bases:

\begin{equation}
v_{n}^{\prime}=\left(\phi_{n}^{\prime}, \Psi_{v}\right)=\left(\sum_{m} S_{m n} \phi_{m}, \sum_{i} v_{i} \phi_{i}\right)=\sum_{m} \sum_{i} S_{m n}^{*} v_{i} \underbrace{\left(\phi_{m}, \phi_{i}\right)}_{\delta_{m i}}=\sum_{m} \sum_{i} S_{m n}^{*} v_{i} \delta_{m i}=\sum_{m} S_{m n}^{*} v_{m}
\end{equation}

Using $S_{m n}^{*}=S_{n m}^{\dagger}$, we can derive the inverse relation:

\begin{equation}
\sum_{i} S_{i n} v_{n}^{\prime}=\sum_{i} S_{i n} \sum_{m} S_{m n}^{*} v_{m}=\sum_{m} \underbrace{\sum_{i} S_{i n} S_{n m}^{\dagger}}_{\delta_{i m}} v_{m}=\sum_{m} \delta_{i m} v_{m}=v_{i}
\end{equation}

Corollary 10.3. Unitary transformations preserve normalization.

Proof. Starting with $\sum_{i}\left|v_{i}\right|^{2}=\left(\Psi_{v}, \Psi_{v}\right)=1$, we can show:

\begin{align}
\left(\Psi_{v}, \Psi_{v}\right) & =\langle v \mid v\rangle=\sum_{m} v_{m}^{*} v_{m}=\sum_{m}\left(\sum_{i} S_{m i} v_{i}^{\prime}\right)^{*}\left(\sum_{j} S_{m j} v_{j}^{\prime}\right)= \\
& =\sum_{m} \sum_{i} \sum_{j} S_{i m}^{\dagger} S_{m j} v_{i}^{\prime*} v_{j}^{\prime}=\sum_{i} \sum_{j} \underbrace{\left(\sum_{m} S_{i m}^{\dagger} S_{m j}\right)}_{\delta_{i j}} v_{i}^{\prime*} v_{j}^{\prime}=  \\
& =\sum_{i} \sum_{j} \delta_{i j} v_{i}^{\prime*} v_{j}^{\prime}=\sum_{i}\left|v_{i}^{\prime}\right|^{2}
\end{align}

\section{Practical implementations}
The Dirac formulation provides elegant representations of quantum dynamics. For the Schrödinger equation:

\begin{equation}
i \hbar \partial_{t} \Psi(x, t)=\hat{H} \Psi(x, t)
\end{equation}

We can express the state vector as:

\begin{equation}
|\Psi\rangle=\sum_{m} c_{m}|m\rangle
\end{equation}

Where $c_{m}=c_{m}(t)$ are time-dependent coefficients. This gives:

\[
\begin{array}{r}
i \hbar \partial_{t}|\Psi\rangle=\hat{H}|\Psi\rangle \\
i \hbar \partial_{t} \sum_{m} c_{m}|m\rangle=\hat{H} \sum_{m} c_{m}|m\rangle
\end{array}
\]

When the basis states are energy eigenstates with $\hat{H}|m\rangle=E_{m}|m\rangle$, the Hamiltonian becomes diagonal:

\begin{equation}
H_{n m}=\langle n| \hat{H}|m\rangle=E_{m}\langle n \mid m\rangle=E_{m} \delta_{n m}
\end{equation}

Resulting in the matrix representation:

\[
H=\left[\begin{array}{llll}
E_{1} & & &  \\
& E_{2} & & \\
& & \ddots & \\
& & & E_{n}
\end{array}\right]
\]

This leads to decoupled differential equations for each coefficient:

\begin{equation}
i \hbar \partial_{t} c_{n}=E_{n} c_{n}
\end{equation}

With solutions:

\begin{equation}
c_{n}(t)=c_{n}(0) \exp \left(-i \frac{t}{\hbar} E_{n}\right)
\end{equation}


And so:

\begin{equation}
|\Psi\rangle=\sum_{m} c_{m}(0) \exp \left(-i \frac{t}{\hbar} E_{m}\right)|m\rangle
\end{equation}

\section{The Heisenberg representation}
In Heisenberg's alternative formulation, we focus on the time evolution of observables rather than states. The expectation value of an operator is:

\begin{equation}
\langle\hat{A}\rangle_{t}=(\Psi(x, t), \hat{A} \Psi(x, t))=\langle\Psi| \tilde{A}|\Psi\rangle
\end{equation}

Introducing the time evolution operator:

\begin{equation}
\hat{U}_{t}=\exp \left(-i \frac{t}{\hbar} \hat{H}\right) \Longrightarrow \Psi(x, t)=\hat{U}_{t} \Psi(x, 0)
\end{equation}

We can express the expectation value in terms of the initial state:

\begin{equation}
(\Psi(x, t), \hat{A} \Psi(x, t))=\left(\hat{U}_{t} \Psi(x, 0), \hat{A} \hat{U}_{t} \Psi(x, 0)\right)=\left(\Psi(x, 0), \hat{U}_{t}^{\dagger} \hat{A} \hat{U}_{t} \Psi(x, 0)\right)
\end{equation}

For notational clarity:

\begin{align}
& \hat{A}=A \\
& \hat{U}_{t}=U  \\
& \Psi(x, 0)=\Psi_{0}
\end{align}

\section{Statistical description with density matrices}
To calculate expectation values, we expand the wavefunction in the energy eigenbasis:

\begin{equation}
\Psi(x, t)=\sum_{r} c_{r} \phi_{r}(x)
\end{equation}

Where $H \phi_{r}(x)=E_{r} \phi_{r}(x)$. The expectation value becomes:

\begin{equation}
\langle A\rangle_{t}=\left(\sum_{r} c_{r} \phi_{r}, A \sum_{s} c_{s} \phi_{s}\right)=\sum_{r} \sum_{s} c_{r}^{*} c_{s}\left(\phi_{r}, A \phi_{s}\right)
\end{equation}

With matrix elements $A_{rs} = (\phi_r, A\phi_s)$ and time-dependent coefficients:

\begin{equation}
c_{r}(t)=c_{r}(0) \exp \left(-i \frac{t}{\hbar} E_{r}\right)
\end{equation}

This yields:

\begin{align}
\langle A\rangle_{t} & =\sum_{r} \sum_{s} c_{r}^{*}(0) \exp \left(i \frac{t}{\hbar} E_{r}\right) c_{s}(0) \exp \left(-i \frac{t}{\hbar} E_{s}\right) A_{r s}= \\
& =\sum_{r} \sum_{s} c_{r}^{*}(0) c_{s}(0) \exp \left(i \frac{t}{\hbar}\left(E_{r}-E_{s}\right)\right) A_{r s}
\end{align}

For sufficiently large times, rapidly oscillating terms with $r \neq s$ average to zero:

\begin{equation}
\langle A\rangle_{t} \simeq \sum_{s}\left|c_{s}(0)\right|^{2} A_{s s}
\end{equation}

In statistical mechanics, the probability distribution follows Boltzmann statistics:

\begin{equation}
\langle A\rangle_{t} \simeq \sum_{s} \frac{\mathrm{e}^{-\beta E_{s}}}{Z} A_{s s}
\end{equation}

These expressions can be written compactly as $\operatorname{Tr}(\rho A)$, where $\rho$ is the density matrix.

\section{Dynamics of observables}
For time-dependent operators in the Heisenberg picture, $A_t = U^\dagger A U$, the equation of motion is:

\begin{equation}
\frac{\mathrm{d} A_{t}}{\, \mathrm{d} t}=\frac{1}{i \hbar}\left[A_{t}, H\right]+\left(\partial_{t} A\right)_{t}
\end{equation}


Proof. Starting from:

\begin{equation}
\frac{\mathrm{d} A_{t}}{\, \mathrm{d} t}=\frac{\mathrm{d}}{\, \mathrm{d} t}\left(U^{\dagger} A U\right)
\end{equation}

The time derivatives of the evolution operators are:

\begin{align}
\frac{\mathrm{d} U}{\, \mathrm{d} t} & =-\frac{i}{\hbar} H U \\
\frac{\mathrm{d} U^{\dagger}}{\, \mathrm{d} t} & =\frac{i}{\hbar} U^{\dagger} H
\end{align}

Applying the product rule:

\begin{align}
\frac{\mathrm{d}}{\, \mathrm{d} t}\left(U^{\dagger} A U\right) & =\frac{\mathrm{d}}{\, \mathrm{d} t}\left(U^{\dagger}\right) A U+U^{\dagger} \frac{\mathrm{d}}{\, \mathrm{d} t}(A) U+U^{\dagger} A \frac{\, \mathrm{d}}{\, \mathrm{d} t}(U)=  \\
& =\frac{i}{\hbar} U^{\dagger} H A U+U^{\dagger} \frac{\, \mathrm{d}}{\, \mathrm{d} t}(A) U-\frac{i}{\hbar} U^{\dagger} A H U
\end{align}

Since $H$ commutes with $U$ (which can be shown using Taylor expansion), we have:

\begin{align}
\frac{\mathrm{d}}{\, \mathrm{d} t}\left(U^{\dagger} A U\right) & =\frac{i}{\hbar} \underbrace{\left(U^{\dagger} H A U-U^{\dagger} A H U\right)}_{\left[U^{\dagger}HU, U^{\dagger}AU\right]}+U^{\dagger} \frac{\, \mathrm{d}}{\, \mathrm{d} t}(A) U= \\
& =\frac{i}{\hbar} U^{\dagger}[H,A]U + U^{\dagger} \frac{\mathrm{d}}{\, \mathrm{d} t}(A) U= \\
& =\frac{1}{i \hbar}\left[A_{t}, H\right]+\left(\partial_{t} A\right)_{t}
\end{align}

This equation closely resembles Ehrenfest's theorem and provides the dynamics of operators in the Heisenberg picture.

\section{Harmonic oscillator example}
For the harmonic oscillator, we define time-evolving operators:

\begin{align}
& x_{t}=U^{\dagger} x U \\
& p_{t}=U^{\dagger} p U
\end{align}

With the Hamiltonian:

\begin{equation}
H\left(p_{t}, x_{t}\right)=\frac{p_{t}^{2}}{2 m}+\frac{1}{2} m \omega^{2} x_{t}^{2}
\end{equation}

Substituting the definitions:

\begin{equation}
H\left(p_{t}, x_{t}\right)=\frac{1}{2 m}\left(U^{\dagger} p U U^{\dagger} p U\right) + \frac{1}{2} m \omega^{2}\left(U^{\dagger} x U U^{\dagger} x U\right)
\end{equation}

Using the unitary property $U U^{\dagger}=\mathbb{I}$:

\begin{equation}
H\left(p_{t}, x_{t}\right)=U^{\dagger}\left(\frac{p^{2}}{2 m}+\frac{1}{2} m \omega^{2} x^{2}\right) U=U^{\dagger} H U=H
\end{equation}

Thus, the Hamiltonian remains time-independent in the Heisenberg picture.

\section{Important unitary transformations}
Several unitary transformations play key roles in quantum mechanics.

Spatial translation operator:

\begin{equation}
T(\lambda)=\exp \left(i \frac{\lambda}{\hbar} p\right)=\exp \left(\lambda \frac{\partial}{\partial x}\right)
\end{equation}

Applying this to a wavefunction yields:

\begin{equation}
T(\lambda) \Psi(x)=\sum_{n} \frac{\lambda^{n}}{n!} \frac{\partial^{n} \Psi(x)}{\partial x^{n}}=\Psi(x+\lambda)
\end{equation}
\section{Rotational transformations}

The angular momentum operator for rotation around the z-axis:

\begin{equation}
L_{3}=-i \hbar\left(x_{1} \frac{\partial}{\partial x_{2}}-x_{2} \frac{\partial}{\partial x_{1}}\right)=-i \hbar \frac{\partial}{\partial \phi}
\end{equation}


The expression in spherical coordinates represents rotation around the z-axis:

\begin{equation}
R(\alpha)=\exp \left(i \frac{\alpha}{\hbar} L_{3}\right)=\exp \left(\alpha \frac{\partial}{\partial \phi}\right)
\end{equation}

This operator rotates functions by angle $\alpha$. For a wavefunction expressed in spherical coordinates:

\begin{equation}
\Psi=\Psi(r \sin (\theta) \cos (\phi), r \sin \theta \sin (\phi), r \cos (\theta))
\end{equation}

The rotation transforms it to:

\begin{equation}
R(\alpha) \Psi=\Psi(r \sin (\theta) \cos (\phi+\alpha), r \sin \theta \sin (\phi+\alpha), r \cos (\theta))
\end{equation}

Using trigonometric addition formulas:

\begin{align}
& r \sin (\theta) \cos (\phi+\alpha)=r \sin (\theta)(\cos (\alpha) \cos (\phi)-\sin (\alpha) \sin (\phi))= \\
& =\underbrace{r \sin (\theta) \cos (\phi)}_{x_{1}} \cos (\alpha)-\underbrace{r \sin (\theta) \sin (\phi)}_{x_{2}} \sin (\alpha)  \\
& r \sin \theta \sin (\phi+\alpha)=r \sin (\theta)(\sin (\alpha) \cos (\phi)+\cos (\alpha) \sin (\phi))= \\
& =\underbrace{r \sin (\theta) \cos (\phi)}_{x_{1}} \sin (\alpha)+\underbrace{r \sin (\theta) \sin (\phi)}_{x_{2}} \cos (\alpha)
\end{align}

This corresponds to the matrix representation:

\[
R(\alpha)=\left[\begin{array}{ccc}
\cos \alpha & -\sin \alpha & 0  \\
\sin \alpha & \cos \alpha & 0 \\
0 & 0 & 1
\end{array}\right]
\]

For time translation, the evolution operator is:

\begin{equation}
U(t+\Delta t)=\exp \left(\frac{-i}{\hbar}(t+\Delta t) H\right)
\end{equation}

Applying this to a wavefunction:

\begin{equation}
U(t+\Delta t) \Psi(x, 0)=\exp \left(\frac{-i}{\hbar} \Delta t H\right) \exp \left(\frac{-i}{\hbar} t H\right) \Psi(x, 0)=\exp \left(\frac{-i}{\hbar} \Delta t H\right) \Psi(x, t)=\Psi(x, t+\Delta t)
\end{equation}

For small time increments $\Delta t \ll t$, using first-order Taylor expansion:

\begin{equation}
U(t+\Delta t) \Psi(x, 0) \sim\left(\mathbb{I}-i \Delta t \frac{H}{\hbar}+\ldots\right) \Psi(x, t)=\Psi(x, t+\Delta t)
\end{equation}

This gives:

\begin{align}
\Psi(x, t+\Delta t) & \sim U(\Delta t) \Psi(x, t)=\mathbb{I} \Psi(x, t)-\Delta t \frac{i}{\hbar} H \Psi(x, t)  \\
\Psi(x, t+\Delta t) & \sim \Psi(x, t)+\frac{\partial \Psi(x, t)}{\partial t} \Delta t
\end{align}

Comparing these expressions:

\begin{align}
\Psi(x, t)-\Delta t \frac{i}{\hbar} H \Psi(x, t) & =\Psi(x, t)+\frac{\partial \Psi(x, t)}{\partial t} \Delta t \\
-\frac{i}{\hbar} H \Psi(x, t) & =\frac{\partial \Psi(x, t)}{\partial t}  \\
H \Psi(x, t) & =i \hbar \frac{\partial \Psi(x, t)}{\partial t}
\end{align}

Thus recovering the Schrödinger equation.
